\documentclass[11pt,a4paper]{amsart}
\pdfinfoomitdate=1
\pdftrailerid{}
\pdfsuppressptexinfo=15
\usepackage[T1]{fontenc}
\usepackage{lmodern}
\usepackage[margin=27mm]{geometry}
\usepackage[expansion=false]{microtype}
\usepackage{amsmath,amssymb,mathtools}
\usepackage{booktabs}
\usepackage{xcolor}
\usepackage[numbers,sort&compress]{natbib}
\usepackage[colorlinks=true,linkcolor=blue!55!black,citecolor=blue!55!black,urlcolor=blue!55!black]{hyperref}
\usepackage[nameinlink,noabbrev]{cleveref}
\raggedbottom

\newtheorem{theorem}{Theorem}[section]
\newtheorem{proposition}[theorem]{Proposition}
\newtheorem{corollary}[theorem]{Corollary}
\newtheorem{lemma}[theorem]{Lemma}
\theoremstyle{definition}
\newtheorem{definition}[theorem]{Definition}
\theoremstyle{remark}
\newtheorem{remark}[theorem]{Remark}

\newcommand{\Cn}{C_n}
\newcommand{\Pin}{\Pi_n}
\newcommand{\Fix}{\operatorname{Fix}}
\newcommand{\card}[1]{\lvert #1\rvert}
\newcommand{\Z}{\mathbb Z}
\newcommand{\Q}{\mathbb Q}

\title[Shift--join partition dynamics]{Cyclic Shift--Join Dynamics on the Partition Lattice:\
M\"obius--Bell Basins and Primitive-Chord Depth}
\author{Anonymous}
\date{Internal Stage 2 draft, 29 August 2026}
\subjclass[2020]{37P35, 05A18, 06B05, 20B25}
\keywords{finite dynamical system, partition lattice, cyclic group, block system, Bell number, M\"obius inversion}

\begin{document}

\begin{abstract}
Let rotation $\rho(i)=i+1$ act on the partition lattice $\Pi_n$ of
$C_n=\mathbb Z/n\mathbb Z$, and iterate
$J_n(\pi)=\pi\vee\rho\pi$.  We analyze this finite functional graph.
The $t$-th iterate is the join of the first $t+1$ translates of $\pi$,
and its endpoint is the coset partition of the subgroup generated by all
within-block differences.  Hence the recurrent points are exactly the
$\tau(n)$ cyclic block systems and
$\zeta_{J_n}(z)=(1-z)^{-\tau(n)}$.  If the endpoint subgroup has order
$h\mid n$, its basin has the exact M\"obius--Bell size
\[
 \sum_{d\mid h}\mu(h/d)B_d^{n/d}.
\]
For $n\geq3$, the sharp stabilization depth is $n-2$.  We also classify its
deepest shell: it consists precisely of partitions with one two-element
block $\{a,a+d\}$ satisfying $\gcd(d,n)=1$ and all remaining blocks
singleton, so its size is $n\varphi(n)/2$.  The classification follows from
a cut-defect argument, not enumeration.  Exact controls make
$1{,}916{,}206$ assertions over all partitions through size ten and,
independently, every binary cut through size twelve.
Classical partition-lattice and permutation block-system results are
owner-subtracted; external release remains on hold.
\end{abstract}

\maketitle

\section{Introduction}

A lattice automorphism and a join define a basic closure dynamics: start at
$x$ and repeatedly adjoin the next translate of the information already
present.  On a general finite lattice, this observation gives convergence
but says little about endpoints, basins, or sharp time.  The partition
lattice over a regular cyclic action is rigid enough to determine all four.

Write $\Pin$ for the partitions of $\Cn=\Z/n\Z$, ordered by refinement,
and let $\rho$ be translation by one.  We study the self-map
\[
             J_n(\pi)=\pi\vee\rho\pi.
\]
The endpoint has a group-theoretic description: take every difference of
two elements lying in one initial block, generate a subgroup, and partition
$\Cn$ into its cosets.  This description makes recurrence elementary and
turns every basin into a divisor-lattice inversion of independent Bell
choices.  Sharp time requires a different argument.  At the penultimate
candidate time, a component cut can have only two untranslated defects.
Parity on the cycles of a chord translation eliminates nonprimitive
differences, while a primitive chord leaves a two-edge boundary that
uniquely recovers that chord.  This forces every deepest initial partition
to be a single primitive-chord atom.

The resulting package has four concrete parts.
\begin{enumerate}
\item We give the exact iterate and the subgroup-coset endpoint of every
  state.
\item We classify recurrence, all iterate-fixed counts, and the finite
  Artin--Mazur zeta function.
\item We count every endpoint basin by M\"obius inversion of Bell powers.
\item We prove the sharp depth $n-2$ and classify all
  $n\varphi(n)/2$ deepest states.
\end{enumerate}

Partition lattices, their joins, Bell enumeration, and ordinary M\"obius
inversion are classical \cite{Stanley2011}.  In permutation-group theory,
invariant partitions are block systems.  Recent work studies lattices of
invariant partitions \cite{AnagnostopoulouMerkouriBaileyCameron2025}, while
Britnell and Wildon study when orbit partitions of group elements are closed
under lattice operations \cite{BritnellWildon2014}.  Those structural facts
including invariant closure by joined translates and the cyclic
coset/subgroup correspondence, receive no credit.  Residual scope is only
the M\"obius--Bell/depth conjunction.  Internally this is neither P97's
subset-sumset map nor P105's permutation-cycle pruning.  A bounded search
through 29 August 2026 did not locate the same package; search absence is not
a novelty or priority certificate.  External dissemination remains
\textbf{HOLD}.

\section{Exact iterates and subgroup endpoints}\label{sec:iterate}

The join $\pi\vee\sigma$ is the least partition coarser than both $\pi$ and
$\sigma$.  Equivalently, its blocks are the connected components of the
graph that joins two vertices whenever they share a block in $\pi$ or in
$\sigma$.  Relabelling by $\rho$ is a lattice automorphism.

\begin{definition}[shift--join map]
For $n\geq1$ and $\pi\in\Pin$, define
\begin{equation}\label{eq:update}
                  J_n(\pi):=\pi\vee\rho\pi .
\end{equation}
The stabilization depth of $\pi$ is
\[
 \tau_n(\pi):=\min\{t\geq0:J_n^{t+1}(\pi)=J_n^t(\pi)\}.
\]
\end{definition}

\begin{proposition}[consecutive-orbit join]\label{prop:iterate}
For every $t\geq0$,
\begin{equation}\label{eq:iterate}
          \boxed{J_n^t(\pi)=\bigvee_{j=0}^{t}\rho^j\pi.}
\end{equation}
\end{proposition}

\begin{proof}
The formula holds at $t=0$.  Suppose it holds at time $t$.  Because $\rho$
preserves joins,
\[
 J_n^{t+1}(\pi)
 =J_n^t(\pi)\vee\rho J_n^t(\pi)
 =\left(\bigvee_{j=0}^{t}\rho^j\pi\right)
  \vee\left(\bigvee_{j=1}^{t+1}\rho^j\pi\right),
\]
which is the join from $j=0$ through $t+1$.
\end{proof}

For a block $B$ of $\pi$, take all differences $x-y$ with $x,y\in B$.
Let
\begin{equation}\label{eq:subgroup}
 H(\pi):=\left\langle x-y:
     x,y\text{ lie in a common block of }\pi\right\rangle\leq\Cn.
\end{equation}
Write $Q_H$ for the partition of $\Cn$ into cosets of a subgroup $H$.

\begin{theorem}[endpoint, recurrence, and zeta]\label{thm:endpoint}
For every $\pi\in\Pin$, the stable endpoint is $Q_{H(\pi)}$.  The fixed and
recurrent states are exactly the coset partitions $Q_H$, one for each
subgroup $H\leq\Cn$.  Consequently, for every $r\geq1$,
\begin{equation}\label{eq:fixed}
 \card{\Fix(J_n^r)}=\tau(n),
\end{equation}
and, as an identity in $\Q[[z]]$,
\begin{equation}\label{eq:zeta}
 \boxed{\displaystyle \zeta_{J_n}(z)
 =\exp\left(\sum_{r\geq1}
   \frac{\card{\Fix(J_n^r)}}{r}z^r\right)
 =(1-z)^{-\tau(n)}.}
\end{equation}
There are $\tau(n)$ one-cycles and no cycles of longer temporal period.
Here $\tau(n)$ denotes the number of positive divisors of $n$.
\end{theorem}

\begin{proof}
The join of all $n$ translates is translation invariant.  It identifies
$0$ with every generator in \eqref{eq:subgroup}, hence with all of
$H(\pi)$, and translation identifies each $x$ with $x+H(\pi)$.  Conversely,
every edge contributed by any translated block has difference in
$H(\pi)$, so it cannot join distinct $H(\pi)$-cosets.  The endpoint is
therefore $Q_{H(\pi)}$.

A fixed partition satisfies $\rho\pi\leq\pi$.  Relabelling preserves the
number of blocks, so this refinement is equality; hence the partition is
invariant under $\rho$ and under the regular action of $\Cn$.  Its block
containing $0$ is a subgroup $H$: if $h\in H$, then
translation by $h$ sends the block containing $0$ to the block containing
$h$, which is the same block; this gives closure under addition, and
finiteness gives inverses.  All other blocks are its cosets.  Conversely,
each $Q_H$ is translation invariant and fixed.  The cyclic group has one
subgroup of each order dividing $n$, so there are $\tau(n)$ fixed points.

Every update coarsens the current partition.  A recurrent orbit in a finite
partially ordered set can therefore contain only equal successive states,
so recurrence equals fixedness.  Thus every positive iterate has the same
fixed set.  Temporal M\"obius inversion gives no cycles of length greater
than one, and the defining series of Artin--Mazur \cite{ArtinMazur1965}
sums to \eqref{eq:zeta}.
\end{proof}

\section{Exact M\"obius--Bell basins}\label{sec:basins}

Let $B_m$ be the $m$-th Bell number.  For $h\mid n$, let
\[
 \mathcal B_n(h):=\{\pi\in\Pin:\card{H(\pi)}=h\}
\]
be the basin of the unique fixed coset partition with block size $h$.

\begin{theorem}[basin formula]\label{thm:basins}
For every $h\mid n$,
\begin{equation}\label{eq:basin}
 \boxed{\displaystyle
 \card{\mathcal B_n(h)}
 =\sum_{d\mid h}\mu(h/d)B_d^{n/d}.}
\end{equation}
In particular, the discrete partition has basin size one, and the one-block
partition has basin
\[
       \sum_{d\mid n}\mu(n/d)B_d^{n/d}.
\]
If $n=p$ is prime, the two basin sizes are $1$ and $B_p-1$.
\end{theorem}

\begin{proof}
Fix the subgroup of order $h$.  The condition $H(\pi)\leq H$ is equivalent
to saying that every block of $\pi$ lies inside one $H$-coset, or that
$\pi$ refines $Q_H$.  There are $n/h$ cosets, each of size $h$, and their
restrictions can be partitioned independently.  Hence
\begin{equation}\label{eq:refinement}
       B_h^{n/h}=\sum_{d\mid h}\card{\mathcal B_n(d)}.
\end{equation}
The subgroup lattice of a cyclic group is the divisor lattice.  Ordinary
M\"obius inversion of \eqref{eq:refinement} proves \eqref{eq:basin}.  The
endpoint cases follow by substitution.
\end{proof}

The basin formula is not inferred from periodic data.  Indeed,
\eqref{eq:zeta} sees only the number $\tau(n)$ of endpoints, whereas
\eqref{eq:basin} retains the full divisor-indexed Bell profile.

\section{Sharp depth and the primitive-chord shell}\label{sec:depth}

For a partition $\pi$, form a graph $G_t(\pi)$ on $\Cn$: for every pair
$a,b$ in a common block of $\pi$ and every $0\leq j\leq t$, add the edge
$\{a+j,b+j\}$.  Its component partition is $J_n^t(\pi)$ by
\cref{prop:iterate}.

\begin{lemma}[one omitted translate is redundant]\label{lem:one-missing}
For every chord $e=\{a,a+d\}$, the graph formed by its first $n-1$
translates has the same connected components as the graph formed by all
$n$ translates.
\end{lemma}

\begin{proof}
Put $g=\gcd(d,n)$ and $\ell=n/g$.  If $\ell\geq3$, the full translate
orbit of $e$ is a disjoint union of $g$ cycles of length $\ell$.  The first
$n-1$ translates omit one edge from one of these cycles, which does not
change its vertex set or connectivity.  If $\ell=2$, each undirected edge
occurs twice among the $n$ translates; omitting one translate leaves every
edge present.  The case $d=0$ contributes no chord.
\end{proof}

Applying the lemma to every within-block pair shows that $G_{n-2}(\pi)$
already has the terminal components.  Thus $\tau_n(\pi)\leq n-2$ for
$n\geq3$.  To determine equality, we need a cut lemma.

For a nonempty proper subset $S\subset\Cn$, write
$f=\mathbf 1_S$.  A chord $e=\{a,a+d\}$ is called \emph{admissible for
$f$} if
\begin{equation}\label{eq:admissible}
       f(a+j)=f(a+d+j),\qquad 0\leq j\leq n-3.
\end{equation}
Thus none of the first $n-2$ translates of $e$ crosses the cut
$S\mid(\Cn\setminus S)$.

\begin{lemma}[two-defect chord lemma]\label{lem:defect}
Assume $n\geq3$ and $f:\Cn\to\{0,1\}$ is nonconstant.  Let
$e=\{a,a+d\}$ be admissible for $f$.
\begin{enumerate}
\item If $\gcd(d,n)>1$, then $f(x)=f(x+d)$ for every $x\in\Cn$.
\item If $\gcd(d,n)=1$, then the only two values of $x$ for which
  $f(x)\ne f(x+d)$ are $a-2$ and $a-1$.  The set $S=f^{-1}(1)$ has no
  nonzero translational stabilizer, and $e$ is the unique primitive chord
  admissible for $f$.
\end{enumerate}
\end{lemma}

\begin{proof}
Set
\[
 D_d(f):=\{x\in\Cn:f(x)\ne f(x+d)\}.
\]
Admissibility gives
$D_d(f)\subseteq\{a-2,a-1\}$.  On every cycle of the permutation
$x\mapsto x+d$, the number of changes of a binary function is even.  If
$\gcd(d,n)>1$, the two possible defects lie in distinct cycles, because
their difference $1$ is not in the subgroup generated by $d$.  Neither
cycle can contain a single defect, so $D_d(f)$ is empty.  This proves the
first assertion.

Suppose now that $d$ is primitive.  Translation by $d$ is one Hamilton
cycle.  The defect set cannot be empty because $f$ is nonconstant, and its
cardinality is even, so
\begin{equation}\label{eq:defects}
             D_d(f)=\{a-2,a-1\}.
\end{equation}
Deleting these two boundary edges splits the Hamilton cycle into two paths;
$f$ is constant on each path and takes different values on them.  Thus $S$
is one nonempty proper cyclic interval in this Hamilton cycle.  A nonzero
rotation cannot preserve a single proper interval: a rotation preserving
its two boundary edges either fixes both, hence is the identity, or swaps
them and exchanges the interval with its complement.  Therefore $S$ has
trivial translational stabilizer.

It remains to prove uniqueness.  Let $s=\card S$, and let
$q\in\{1,\ldots,n-1\}$ satisfy $qd\equiv1\pmod n$.  The starting vertices
$a-2$ and $a-1$ of the two boundary edges are separated by $q$ steps in the
$d$-cycle.  The two components after deleting them have $q$ and $n-q$
vertices.  Hence $s\in\{q,n-q\}$ and
\begin{equation}\label{eq:recover-d}
                   d\equiv\pm s^{-1}\pmod n.
\end{equation}
Any other admissible primitive chord must satisfy the same equation, so it
has the same undirected difference and the same Hamilton cycle.  Its two
omitted translates must be exactly the two cut edges determined by $S$.
After orienting the difference as $d$, write those consecutive edges as
$E_p=\{p,p+d\}$ and $E_{p+1}$.  The initial chord must be $E_{p+2}$;
reversing the orientation of $d$ gives the same undirected edge.  Thus the
initial chord is uniquely $e$.
\end{proof}

\begin{theorem}[sharp depth and complete deepest shell]\label{thm:depth}
For $n\geq3$,
\begin{equation}\label{eq:maxdepth}
        \max_{\pi\in\Pin}\tau_n(\pi)=n-2.
\end{equation}
Moreover, $\tau_n(\pi)=n-2$ if and only if $\pi$ has one two-element block
\begin{equation}\label{eq:primitive-atom}
        \{a,a+d\},\qquad \gcd(d,n)=1,
\end{equation}
and all other blocks are singleton.  Consequently,
\begin{equation}\label{eq:deepest-count}
 \boxed{\displaystyle
 \#\{\pi\in\Pin:\tau_n(\pi)=n-2\}
 =\frac{n\varphi(n)}2.}
\end{equation}
For $n=1,2$, every partition is fixed and the maximum depth is zero.
\end{theorem}

\begin{proof}
The upper bound in \eqref{eq:maxdepth} follows from
\cref{lem:one-missing}.  If $\pi$ is the atom in
\eqref{eq:primitive-atom}, its full translate graph is a Hamilton cycle.
At time $n-3$, exactly two distinct edges of that cycle are absent, so the
graph is disconnected.  At time $n-2$, \cref{lem:one-missing} gives the
full connected component.  Hence this atom has depth $n-2$.

Conversely, suppose $\tau_n(\pi)=n-2$.  Some component $S$ of
$G_{n-3}(\pi)$ is a nonempty proper subset of its component in the terminal
translate graph.  Let $f=\mathbf1_S$.  Every chord joining two elements of
one initial block is admissible for $f$, since all its first $n-2$
translates are edges of $G_{n-3}(\pi)$ and no such edge crosses a component
cut.

If every initial chord had nonprimitive difference, the first part of
\cref{lem:defect} would make $f$ invariant under every within-block
difference.  It would then be invariant under the subgroup $H(\pi)$ they
generate.  Thus $S$ would be a union of entire $H(\pi)$-cosets, contrary to
its being a proper nonempty subset of one terminal coset.  There is
therefore an admissible primitive chord $e$.  By the second part of
\cref{lem:defect}, $S$ has trivial translational stabilizer and $e$ is the
only admissible primitive chord.  A second nonprimitive chord would give a
nonzero translational stabilizer by the first part of the lemma, while a
second primitive chord would violate uniqueness.  Hence $e$ is the only
chord in the equivalence graph of $\pi$.  Since every block contributes all
of its internal chords, $\pi$ has exactly the one two-element block $e$ and
all remaining blocks are singleton.

Finally, choose $a\in\Cn$ and a unit $d\in\Cn$.  This gives
$n\varphi(n)$ oriented primitive chords.  Each undirected chord is counted
twice, by $(a,d)$ and $(a+d,-d)$, proving
\eqref{eq:deepest-count}.  For $n=1,2$, direct inspection gives the stated
endpoint convention.
\end{proof}

The proof of \cref{thm:depth} is independent of the M\"obius--Bell basin
route.  Its only enumerative step is the final count of primitive chords;
the if-and-only-if classification is a graph-cut theorem.

\section{Recovery, controls, and limitations}\label{sec:control}

The periodic ledger is deliberately incomplete as a parameter invariant:
different values of $n$ with the same divisor count have the same zeta
function.  Transient time restores recovery for $n\geq3$, because
$n=\max\tau_n+2$.  The two small cases have the same maximum depth, but
their phase sizes $B_1=1$ and $B_2=2$ distinguish them.

Two analytic routes support the result.  The \emph{semilattice--subgroup
route} proves the iterate, endpoint, recurrence, zeta, and basin formula.
The \emph{translated-graph route} realizes joins as components and proves
sharp depth through chord cycles and binary cut defects.  Neither route is
a numerical interpolation of the other.

The deterministic control in \texttt{code/verify.py} supplies a third,
finite check.  It enumerates every restricted-growth string through
$n=10$, compares literal shift--join orbits with both a separately
accumulated lattice join and an independently built translated graph,
checks subgroup endpoints, all divisor-indexed basins, and the deepest
classification state by state.  It independently exhausts every binary cut
through $n=12$ for the defect lemma, the basin convolution through $n=50$,
and temporal M\"obius/zeta identities through period $60$.  The fresh run
contains $142{,}417$ literal partitions and $1{,}916{,}206$ exact integer
assertions, with no randomness or floating-point calculation.

The results do not count every intermediate depth layer or every one-step
fiber, and search absence does not establish priority.  Residual scope is
limited to the basin/depth conjunction proved here.  Public release and
specialist contact remain \textbf{HOLD}.

\section{Conclusion}

The map has exact subgroup endpoints and M\"obius--Bell basins.  Its deepest
shell consists of primitive chords, has size $n\varphi(n)/2$, and occurs at
sharp depth $n-2$.  Extending the basin and sharp-depth analysis to noncyclic
regular actions is a natural next problem.

\begingroup
\footnotesize
\bibliographystyle{plainnat}
\bibliography{references}
\endgroup

\end{document}
